\documentclass{article}
\usepackage{amsmath, amssymb, amsthm}
\usepackage{hyperref}
\usepackage{geometry}
\geometry{margin=1in}

\title{Erd\H{o}s Problem \#876}
\date{Last updated: 2025-08-31}
\author{Source: erdosproblems.com}

\begin{document}
\maketitle

\noindent\textbf{Status:} OPEN \quad \textbf{No prize}

\noindent\textbf{Tags:} additive combinatorics

\medskip

\noindent\textbf{Problem Statement:}

Let $A=\{a_1&#60;a_2&#60;\cdots\}\subset \mathbb{N}$ be an infinite sum-free set - that is, there are no solutions to\[a=b_1+\cdots+b_r\]with $b_1&#60;\cdots&#60;b_r&#60;a\in A$. How small can $a_{n+1}-a_n$ be? Is it possible that $a_{n+1}-a_n&#60;n$?




Erd\H{o}s \cite{Er98} writes that Graham 'recently proved' that there is such a sequence for which $a_{n+1}-a_n&lt;n^{1+o(1)}$, and that Melfi proved a somewhat weaker result. Erd\H{o}s \cite{Er62c} proved that a sum-free set has density zero. Deshouillers, Erd\H{o}s, and Melfi \cite{DEM99} constructed a sum-free set that grows like $a_n\sim n^{3+o(1)}$. Luczak and Schoen \cite{LuSc00} have proved that, for all large $N$,\[\lvert A\cap [1,N]\rvert\ll (N\log N)^{1/2},\]and that there exists a sum-free set $B$ such that\[\lvert B\cap [1,N]\rvert \gg \frac{N^{1/2}}{(\log N)^{1/2+o(1)}}\]for all large $N$. In \cite{Er75b} and \cite{Er77c} Erd\H{o}s asks to determine the maximum possible value of $\sum_{n\in A}\frac{1}{n}$. Erd\H{o}s had proved this is $&lt;100$, and Sullivan had shown that this is $&lt;4$, and Sullivan conjectured the maximum is slightly larger than $2$. See also [790] .




References

[DEM99] Deshouillers, Jean-Marc and Erd\H{o}s, Paul and Melfi, Giuseppe, On a question about sum-free sequences . Discrete Math. (1999), 49--54.

[Er62c] Erd\H{o}s, P\'{a}l, Some remarks on number theory. {III} . Mat. Lapok (1962), 28--38.

[Er75b] Erd\H{o}s, Paul, Problems and results in combinatorial number theory . Journ\'{e}es Arithm\'{e}tiques de Bordeaux (Conf., Univ. Bordeaux, Bordeaux, 1974) (1975), 295-310.

[Er77c] Erd\H{o}s, Paul, Problems and results on combinatorial number theory. III . Number theory day (Proc. Conf., Rockefeller Univ., New York, 1976) (1977), 43-72.

[Er98] Erd\H{o}s, Paul, Some of my new and almost new problems and results in combinatorial number theory . Number theory (Eger, 1996) (1998), 169-180.

[LuSc00] \L uczak, Tomasz and Schoen, Tomasz, On the maximal density of sum-free sets . Acta Arith. (2000), 225--229.


\bigskip
\noindent\small{Source: \url{https://www.erdosproblems.com/876}}
\end{document}
